\section{Extensions and field Morita finiteness}\label{sec:extensions}

Throughout this section, a bound is allowed to depend on the fixed prime
$p$ and the defect-order bound $M$, but on no ambient group order.  Put
$k=\overline{\F}_p$ and $\cO=W(k)$, and let $\sigma$ denote coefficient
Frobenius.  We use integral blocks at intermediate stages.  The conclusion
for arbitrary finite groups will be a conclusion over $k$.

The input from the preceding sections is the quasisimple Cartan
bound.  We first reduce a general block to a crossed extension
whose identity component has bounded defect and Cartan entries.
We then specify an equivariant Frobenius comparison and prove
that it makes these extensions finite up to Morita equivalence.
Section~\ref{sec:periodicity} will construct that comparison
without using the finite lists obtained here.

\subsection{Reduction to controlled crossed extensions}

Crossed products and algebraic-group actions enter Donovan's
conjecture through K\"ulshammer's reduction
\cite{Kulshammer1995}; Eisele developed its integral form
\cite{EiseleReduction}.  We retain both the outer action and the
unit-valued multiplication factors.  The additional issue here is
comparison through Sylow $p$-actions, beyond the prime-to-$p$
crossed-product finiteness of those reductions.

\begin{lemma}[Structure of a reduced pair]\label{lem:ext:reduced-structure}
Every block of a finite group with defect order at most $M$ is
$k$-linearly Morita equivalent to a block $B$ of a finite group $H$ with
the following properties.
\begin{enumerate}
\item Every block of a normal subgroup covered by $B$ is $H$-stable.
If $B$ covers a nilpotent block of $H_0\trianglelefteq H$, then
$H_0\leq Z(H)O_p(H)$.
\item Write
\[
 \begin{gathered}
 N=F^*(H)=PZE(H),\qquad E(H)=K_1\cdots K_j,\\
 P=O_p(H),\qquad Z=O_{p'}(H),
 \end{gathered}
\]
where the $K_i$ are the components, that is, the subnormal quasisimple
subgroups.  Then $Z\leq Z(H)$, $|P|\leq M$, and $j$ is bounded.  The
block $b$ of $N$ covered by $B$ has bounded defect and bounded Cartan
entries.
\item For $X=H/N$, the index $[X:O_{p'}(X)]$ is bounded.  The same
index bound holds, with a possibly different constant, for every
subgroup of every extension of $X$ by a $p'$-group.
\end{enumerate}
\end{lemma}

\begin{proof}
The general reduction in An--Eaton, Proposition~6.1
\cite{AnEaton}, gives the first assertion and preserves the defect
group.  Its statement has no restriction on the defect group.  Only
this preliminary reduction is used here.  If necessary it can be
performed over a sufficiently large modular system and then reduced
to $k$.  All subsequent integral blocks are the canonical lifts of
blocks of the resulting finite groups to $\cO$.

A block of the normal $p'$-subgroup $O_{p'}(H)$ is nilpotent, so the
first assertion makes that subgroup central.  Also $P$ is contained
in a defect group of $B$.  The displayed description of $N$ follows
from the definition of the generalized Fitting subgroup.  Its factors
commute except for their central intersections.

Let $b_i$ be a block of $K_i$ covered by $b$, with defect group $D_i$.
Normal block theory, applied also along a subnormal chain, gives
$|D_i|\leq M$.  No $D_i$ is central in $K_i$.  Indeed, all components
in the $H$-orbit of such a component would have central-defect blocks.
Their central product is normal in $H$, and its covered block has
central defect and is therefore nilpotent.  The first assertion would
put this nontrivial perfect group inside the soluble group $Z(H)P$,
which is impossible.

The multiplication map
\[
 P\times Z\times\prod_{i=1}^j K_i\longrightarrow N
\]
has central kernel.  Lift $b$ along this map, taking the trivial
character on the $p'$-part of the kernel.  Its lift is a tensor product
of the block of $P$, a linear-character block of $Z$, and the $b_i$.
The kernel identifies only central elements.  Consequently the
noncentral quotients $D_i/(D_i\cap Z(K_i))$ contribute independently
to a defect group of $b$.  Each has order at least $p$, and hence
$p^j\leq M$.  The lifted tensor block has defect order at most
$M^{j+1}$.  Theorem~\ref{thm:quasisimple-cartan}, the bound on $|P|$,
and the tensor-product formula bound its Cartan entries.  Descent
through the central kernel, using Lemma~\ref{lem:central-transfer},
then bounds the Cartan entries of $b$.

We next bound the quotient.  The generalized Fitting centralizer
theorem \cite[(31.13)]{Aschbacher} says that
$C_H(F^*(H))\leq F^*(H)$.  It embeds $X$ into the
group of outer actions on $P$ and the permuting outer actions on the
components.  To check the kernel, inner actions on $P$ and on each
component can be removed simultaneously by elements of their
commuting factors in $N$.  An element left in the kernel centralizes
$N$, since it also centralizes $Z$, and therefore belongs to $N$.
The standard outer automorphism descriptions of finite simple groups
give a normal soluble subgroup $X_0$ of bounded index and bounded
derived length in $X$: the component permutation group and $\Out(P)$
have bounded order.  For Lie types, the diagonal group is abelian,
the field group is cyclic, and the diagram group has bounded
derived length \cite[Theorem~3.4]{BMO}; for alternating groups the
outer group is of order at most two outside the single degree-six
exception \cite[Theorem~2.3 and Section~2.4.2]{Wilson}; and the sporadic list is
finite.  Thus the outer groups have uniformly bounded derived
length.  The same statements hold for their perfect
central covers, since an automorphism trivial on the simple quotient
and on inner automorphisms is trivial on the perfect cover.

Consider an abelian derived section $X_0^{(i)}/X_0^{(i+1)}$.  Its
$p$-primary subgroup lifts to a normal section $H_2/H_1$ of $H$,
above $N$, of $p$-power order.  All the relevant covered blocks are
stable.  The defect group of the covered block of $H_2$ surjects onto
$H_2/H_1$, by the normal block theorem for a $p$-power extension.
It has order at most $M$.  Thus the $p$-part of each derived section
has order at most $M$, and $|X|_p$ is bounded.

Set $Q=X/O_{p'}(X)$.  Then $O_{p'}(Q)=1$, so $F(Q)=O_p(Q)$ has
bounded order.  All nonabelian composition factors of $Q$, counted
with multiplicity, occur in a quotient of bounded order, namely
$X/X_0$; their number and orders are bounded.  Components of $Q$
are quotients of the universal central covers of these finitely many
simple groups.  Therefore $F^*(Q)$ has bounded order.  Applying the
generalized Fitting centralizer theorem to $Q$ shows that $Q$ has
bounded order as well: its conjugation homomorphism into
$\Aut(F^*(Q))$ has kernel contained in $F^*(Q)$.

Finally, if $\widetilde X\twoheadrightarrow X$ has $p'$-kernel, the
inverse image of $O_{p'}(X)$ is a normal $p'$-subgroup of bounded
index.  Its intersection with any subgroup of $\widetilde X$ proves
the last assertion.
\end{proof}

We next enlarge the components in a way which preserves both the
defect bound downstairs and the actual multiplication of the outer
actions.  The second requirement is essential for the later crossed
products.

\begin{lemma}[Compatible partial regular extensions]
\label{lem:ext:partial-extension}
In the notation of Lemma~\ref{lem:ext:reduced-structure}, there is an
extension $N\trianglelefteq\bar N$ with abelian $p'$-quotient
$A=\bar N/N$ having the following properties.
\begin{enumerate}
\item Choose $h_1=1$ and representatives $h_x$ for $x\in X$, and put
\begin{equation}\label{eq:ext:actual-factors}
 h_xh_y=n_{xy}h_{xy},\qquad n_{xy}\in N.
\end{equation}
Conjugation by $h_x$ on $N$ extends to an automorphism $\alpha_x$ of
$\bar N$, and
\begin{align}
 \alpha_x\alpha_y&=\operatorname{ad}(n_{xy})\alpha_{xy},
 \label{eq:ext:action-law}\\
 n_{xy}n_{xy,z}&=\alpha_x(n_{yz})n_{x,yz}.
 \label{eq:ext:factor-law}
\end{align}
Here $\operatorname{ad}(g)$ means conjugation by $g$; the factors are
the actual elements in \eqref{eq:ext:actual-factors}.
\item Every block $\bar b$ of $\bar N$ covering $b$ has bounded defect
and Cartan entries.
\item There is a central presentation
\[
 \bar N=\bigl(P\times Z\times\prod_i V_i\bigr)/D_0,
 \qquad
 N=\bigl(P\times Z\times\prod_i U_i\bigr)/D_0,
\]
where $U_i$ is the universal cover of the simple quotient of $K_i$.
Outside alternating types and a finite list, $U_i=\mathbf U_i^{F_i}$
is a standard simply connected group in characteristic different
from $p$.  It has a reductive enlargement with fixed-point group
\[
 J_i=V_i\times C_i,
\]
where $C_i$ is central of $p$-power order, $V_i/U_i$ is abelian of
$p'$-order, the derived algebraic group is simply connected, and
all dual semisimple $p'$-centralizers are connected.  The constructions
respect the component permutations and the actions in the first
assertion.  No bound on $|D_0|$ or $|C_i|$ is asserted.
\end{enumerate}
\end{lemma}

\begin{proof}
The universal central extension is functorial for perfect groups.
It therefore lifts the actions on the components and gives a central
surjection $P\times Z\times\prod_iU_i\to N$ with kernel $D_0$.
In defining characteristic, a component block with noncentral defect
has a Sylow defect group, apart from the bounded exceptional groups;
see the defining-characteristic block theorem \cite{HumphreysDefining}.
A bound on this Sylow order bounds both field degree and rank.
Thus these components, the sporadic components, and the exceptional
multiplier and automorphism cases form a finite list.  The remaining
nonalternating factors are the standard simply connected finite
groups in cross characteristic.

For one such factor write $U=\mathbf U^F$, with a pinned simply
connected simple algebraic group $\mathbf U$.  First suppose the
Frobenius has the usual form $F=\gamma\operatorname{Fr}_r^f$, with
$\gamma$ a diagram automorphism and $r\ne p$.  Let
$Z_{p'}=Z(\mathbf U)_{p'}$ be the prime-to-$p$ part of its finite
diagonalizable center, interpreted as a group scheme.  Take a split
torus $\mathbf T_0$ with one coordinate for each fundamental weight.
Restriction of these weights to $Z_{p'}$ embeds $Z_{p'}$ into
$\mathbf T_0$.  This embedding is equivariant for diagram
permutations and split Frobenius.  Define
\begin{equation}\label{eq:ext:partial-regular}
 \mathbf J=(\mathbf U\times\mathbf T_0)/
 \{(z,z^{-1}):z\in Z_{p'}\}.
\end{equation}
The construction includes nonreduced diagonalizable parts when the
defining characteristic divides the center order.  In particular it
is a construction on the pinned root datum, rather than merely on
abstract finite centers.

The original $\mathbf U$ embeds as $[\mathbf J,\mathbf J]$, and is
still simply connected.  As group schemes its center satisfies
\[
 Z(\mathbf J)\simeq Z(\mathbf U)_p\times\mathbf T_0.
\]
The center component group is therefore $p$-primary.  The obstruction
in the fundamental group to
connected semisimple centralizers on the dual side is $p$-primary.
The semisimple centralizer component theorem embeds the component
group in that obstruction and makes its exponent divide the order
of the semisimple element; see
\cite[Corollary~2.9]{BonnafeQuasiIsolated} for connected reductive
groups and \cite[Proposition~14.20]{MalleTesterman}.
It follows that a semisimple $p'$-element of $\mathbf J^*$ has
connected centralizer.

Put $J=\mathbf J^F$.  Lang--Steinberg's theorem
\cite[Theorem~21.7]{MalleTesterman} applied to the connected kernel
$\mathbf U$ identifies
\[
 J/U=(\mathbf T_0/Z_{p'})^F.
\]
The isogeny $\mathbf T_0\to\mathbf T_0/Z_{p'}$ has degree prime to
$p$.  On fixed points its kernel and cokernel have $p'$-order; for
the cokernel use the fixed-point exact sequence and $H^1(F,Z_{p'})$.
It therefore induces an isomorphism on $p$-primary subgroups.  The
$p$-primary subgroup of $\mathbf T_0^F$ gives a central subgroup
$C\leq J$ mapping isomorphically onto $(J/U)_p$ and disjoint from
$U$.  Let $V$ be the inverse image of $(J/U)_{p'}$.  Then
$J=V\times C$ and $V/U$ is abelian of $p'$-order.  These definitions
are invariant under every automorphism of the pinned construction.

For exceptional graph isogenies in types $B_2,F_4$ in characteristic
$2$, and $G_2$ in characteristic $3$, use the original group on
points and put $V=U$, $C=1$.  This includes the Suzuki and Ree
endomorphisms.  The center and the relevant fundamental obstruction
have trivial groups of geometric points in the only nontrivial
isogeny case, or are trivial; the same connected-centralizer
conclusion holds for semisimple elements.  For the finite-list and
alternating factors also put $V_i=U_i$.

We spell out compatibility of automorphisms.  For a standard
universal Lie group, the automorphism theorem
\cite[Definition~3.3 and Theorem~3.4]{BMO} gives the semidirect
product of the rational inner-diagonal group and the concrete
field-and-graph subgroup.  In the ordinary diagram case, write
$\mathcal E=\langle\operatorname{Fr}_r,\text{pinned diagrams}\rangle$
on geometric points.  The restriction of $C_{\mathcal E}(F)$ to
$U$ has kernel $\langle F\rangle$; see \cite[Section~2.2]{SpaethD}.
Thus its image is exactly the pinned field-and-diagram group with
that relation imposed.  The inner-diagonal group is represented
faithfully by $\mathbf U_{\rm ad}^F$.  It acts on
\eqref{eq:ext:partial-regular} by conjugation on $\mathbf U$ and
trivially on the torus.  The pinned operations act on $\mathbf U$
and on the fundamental-weight coordinates of $\mathbf T_0$ with
the same relations and the same conjugation action on the
inner-diagonal group.  The sole additional relation
$F=\gamma\operatorname{Fr}_r^f=1$ holds on $J$.  Consequently the
semidirect product action factors through the actual automorphism
group of $U$.  An inner automorphism implemented by $u\in U$ is
extended by conjugation by that same element.  Use identical
pinned models on isomorphic permuted components.  In the
exceptional graph-isogeny case no enlargement is made, so the
existing actions already have the required property.

Apply this construction factor by factor.  The old kernel $D_0$
remains central and is preserved, since the extended operations
agree with the old operations on all its elements.  Taking its
quotient gives $\bar N$ and an abelian $p'$-quotient $A$ over $N$.
The compatibility just proved makes the products of the extended
actions differ by conjugation by exactly $n_{xy}$, proving
\eqref{eq:ext:action-law}.  Equation~\eqref{eq:ext:factor-law}
already holds as an equality in $N$, by associativity in $H$, and
therefore still holds in $\bar N$.

Finally, $N\trianglelefteq\bar N$ has $p'$-index, so a block covering
$b$ has the same defect order.  Lemma~\ref{lem:abelian-transfer}
bounds its Cartan entries.  This argument takes place after
quotienting by $D_0$; it does not require bounded defect for a block
lifted to the universal covers or to the groups $J_i$.
\end{proof}

\begin{lemma}[Recovery by a dual action]\label{lem:ext:dual-recovery}
Let $\Xi=\Hom(A,k^\times)$ and $Y=\Xi\rtimes X$.  There is an
integral crossed algebra with identity component $\cO\bar N$ and
grading group $Y$ such that $\cO H$ is a full idempotent corner.
After taking the summand relevant to $B$, and then a further full
corner, one obtains a crossed algebra with identity component
$\bar B=\cO\bar N\bar b$ and grading group
$Y_0=\operatorname{Stab}_Y(\bar b)$.  Every such $\bar b$ covers $b$.  The groups
$Y_0/O_{p'}(Y_0)$, and hence their Sylow $p$-subgroups, have bounded
order.
\end{lemma}

\begin{proof}
For $\chi\in\Xi$, let its action on a group element $g\in\bar N$
be multiplication by $\chi(gN)$.  Lift these characters by
Teichm\"uller representatives over $\cO$.  Combine these actions
with the $\alpha_x$ of Lemma~\ref{lem:ext:partial-extension}; use
$n_{xy}$ as the factor for the $X$-coordinates and ordinary
semidirect transport on $\Xi$.  Equations~\eqref{eq:ext:action-law}
and \eqref{eq:ext:factor-law} give a crossed system, since every
character in $\Xi$ is trivial on all $n_{xy}\in N$.

Write $u_\chi$ for the homogeneous units corresponding to $\Xi$.
The element
\[
 e_\Xi=\frac1{|\Xi|}\sum_{\chi\in\Xi}u_\chi
\]
is an idempotent.  A group element in the $a$-component of the
$A$-grading of $\cO\bar N$ conjugates $e_\Xi$ to the character
idempotent of $\cO\Xi$ corresponding to $a$.  These idempotents,
as $a$ varies, sum to $1$, so $e_\Xi$ is full.  Moreover
$e_\Xi(\cO\bar N\rtimes\Xi)e_\Xi=\cO N e_\Xi$:
the average kills all nonidentity $A$-components.  The $X$-units
commute with $e_\Xi$ and retain their factors $n_{xy}$.  Its corner
in the entire crossed algebra is consequently $\cO H$.

Decompose the crossed algebra by $Y$-orbits on the blocks of
$\cO\bar N$.  Each orbit sum is fixed by $\Xi$ and so belongs to
$\cO N$.  The orbit summand meeting $B$ has nonzero product with
$b$.  At least one of its blocks covers $b$; every member does,
since twists are trivial on $N$ and the $H$-actions preserve $b$.
Within this orbit summand a single block idempotent $\bar b$ is
full, since its homogeneous conjugates sum to the orbit idempotent.
Its corner has precisely the degrees in $Y_0$, with invertible
homogeneous units $\bar b u_y$, and is the asserted crossed
algebra.  These constructions also hold after reduction to $k$.

The group $Y$ is an extension of $X$ by the $p'$-group $\Xi$.
The bound follows from Lemma~\ref{lem:ext:reduced-structure}.
\end{proof}

\subsection{The comparison input and integral base orders}

The reduction has bounded the defect and Cartan data of the
identity blocks, and bounded the grading groups modulo their
normal $p'$-subgroups.  The grading itself may still be large.
We now state the additional comparison needed to control its
actual crossed multiplication.

Here and below a crossed system on an algebra $R$ for a finite group
$T$ consists of automorphisms $\alpha_t$ and units $a_{s,t}$ such that
\begin{align*}
 \alpha_s\alpha_t&=\operatorname{ad}(a_{s,t})\alpha_{st},\\
 a_{s,t}a_{st,u}&=\alpha_s(a_{t,u})a_{s,tu},
\end{align*}
with the usual identity normalizations.  Its algebra is
$\bigoplus_{t\in T}Ru_t$, with $u_t r=\alpha_t(r)u_t$ and
$u_su_t=a_{s,t}u_{st}$.  An isomorphism is called \emph{based} if it
is the identity on the specified copy of $R$ and preserves each
group-labelled homogeneous component.  Such isomorphisms are
exactly the changes of homogeneous units.
For the crossed-system formalism, see also
\cite[Definition~4.1 and Lemma~4.3]{EiseleReduction}.

\begin{definition}[The comparison property]\label{def:ext:comparison-data}
For the data of Lemmas~\ref{lem:ext:partial-extension} and
\ref{lem:ext:dual-recovery}, let $S\leq Y_0$ be a $p$-subgroup
contained in the pure subgroup $X\leq\Xi\rtimes X$.  The comparison
property is the existence of a positive integer $m$, bounded in
terms of $p,M$, and an integral
$(\sigma^m\bar B,\bar B)$-Morita bimodule $\mathcal M$ with the
following additional structure on $M_0=k\otimes_{\cO}\mathcal M$.
There are invertible $k$-linear maps $J_x:M_0\to M_0$, $x\in S$,
such that, writing $\alpha'_x=\sigma^m(\alpha_x)$ and
$a'_{xy}=\sigma^m(a_{xy})$, one has
\begin{align}
 J_x(avb)&=\alpha'_x(a)J_x(v)\alpha_x(b),
 \label{eq:ext:comparison-covariance}\\
 J_xJ_y(v)&=a'_{xy}J_{xy}(v)a_{xy}^{-1},\qquad J_1=\id.
 \label{eq:ext:comparison-law}
\end{align}
Here $a_{xy}=n_{xy}\bar b$ and the target factor is its coefficient
Frobenius image, with the corresponding target block identity.
The property is required also for $S=1$.
\end{definition}

It is enough to produce these maps with a scalar error in
\eqref{eq:ext:comparison-law}.  Associativity and covariance then
make the errors a scalar $2$-cocycle on $S$.  Positive cohomology
of $S$ with coefficients in $k^\times$ vanishes: $|S|$ annihilates
it, whereas the $|S|$-power map of $k^\times$ is an automorphism.
Rescaling the $J_x$ removes the error.  An error in arbitrary units
of the center would not give this conclusion.

Every $p$-subgroup of $Y$ is conjugate by an element of $\Xi$ to a
subgroup of the pure $X$.  Indeed its image in $X$ is a $p$-group,
and Schur--Zassenhaus conjugates its complement to the standard
complement in the inverse image of that image.  Conjugating also
the block and the corner transports this assertion to every
$p$-subgroup of $Y_0$.  This conjugation is implemented by an
integral homogeneous unit in the algebra of
Lemma~\ref{lem:ext:dual-recovery}.

\begin{lemma}[Finite integral bases and finite Picard images]
\label{lem:ext:integral-bases}
Assume the comparison property for the trivial subgroup.  The
integral basic orders of all the blocks $\bar B$ above belong to a
finite list $\mathfrak R_1,\ldots,\mathfrak R_t$.  Each
$\mathfrak R_i$ is defined over some $W(\F_{p^{c_i}})$ and has finite
$\Pic_{\cO}(\mathfrak R_i)$.  Consequently the outer actions in
the reduced crossed basic corners take their values in a fixed
finite subgroup of $\Pic_k(R_i)$, where
$R_i=k\otimes_{\cO}\mathfrak R_i$.
\end{lemma}

\begin{proof}
The defect orders and Cartan entries of $\bar B$ are bounded.
The number of simple modules is bounded by the defect bound, so
the sum of its Cartan entries is bounded too.  For $S=1$ the
comparison property bounds the integral Morita--Frobenius number.
Eaton--Eisele--Livesey, Theorem~3.10 \cite{EEL}, says precisely
that bounded Cartan sum, bounded numerical defect $d$ (so the
defect group has order $p^d$), and bounded integral
Morita--Frobenius number give finitely many integral Morita classes.
Applying it to the finitely many possible numerical defects gives
the claimed finite list of basic orders.

Choose each representative as a basic corner of one actual
integral group block.  Its block and primitive idempotents modulo
$p$ are defined over a finite field containing their finitely many
coefficients and splitting the finitely many simple modules in
question.  Lift these idempotents over its Witt ring, which is
complete.  The resulting corner defines $\mathfrak R_i$ over
$W(\F_{p^{c_i}})$ for some $c_i$.

Eisele's Theorem~B and Corollary~1.2 \cite{EiselePicard} give the
finiteness of the integral Picard group of a block, and hence of
its basic order.  The hypotheses apply to $\cO=W(k)$: it is
unramified and has algebraically closed residue field, and the
generic algebra is separable.  The relevant rigidity condition can
also be checked directly.  The conjugacy-class decomposition of
the adjoint group module gives
\[
 \operatorname{HH}^1(\cO G)
   =\bigoplus_{[g]}H^1(C_G(g),\cO)=0,
\]
because each centralizer is finite and $\cO$ is torsion-free with
trivial action in the displayed cohomology.  Vanishing passes to
block summands and to basic orders by Morita invariance, as
required in Eisele's theorem.

Passing to basic corners preserves a crossed structure integrally;
see \cite[Proposition~4.15 and Corollary~4.16(2)]{EiseleReduction}.
We recall how the homogeneous units and their factors are obtained.
If $e$ is a basic idempotent in a block order, then any automorphism
sends $e$ to a unit-conjugate
idempotent: both associated projectives contain one copy of every
indecomposable projective type.  Multiplying each homogeneous unit
by a suitable identity-component unit makes it commute with $e$.
The corner is then crossed over $e\bar B e$.  After identifying
this order with $\mathfrak R_i$, its homogeneous automorphisms
give elements of $\Pic_{\cO}(\mathfrak R_i)$.  Their reductions
lie in the finite image of this group in $\Pic_k(R_i)$.  Since
$R_i$ is basic, $\Pic_k(R_i)$ is naturally its outer automorphism
group: an invertible bimodule is isomorphic to $R_i$ as a one-sided
module, and its other action specifies an automorphism up to
inner automorphism.
\end{proof}

\subsection{Based descent and the large kernel}

The trivial-subgroup comparison has supplied finitely many
integral base orders and finite images of their Picard groups.
We next use the full Sylow comparison to retain the chosen base
algebra while descending the crossed factors.  This stronger
form of finiteness is what permits removal of a large $p'$-kernel.

\begin{lemma}[Based Frobenius descent on $p$-subgroups]
\label{lem:ext:based-sylow}
Assume the comparison property of
Definition~\ref{def:ext:comparison-data}.  For each fixed integral
basic order $\mathfrak R_i$ and each group of bounded $p$-power
order, the restrictions of the reduced crossed basic corners to
that group have only finitely many based isomorphism classes.
The identification of the identity component with $R_i$ is part
of this assertion and may be chosen arbitrarily integrally.
\end{lemma}

\begin{proof}
The argument has three steps.  We align the chosen basic corners
with their Frobenius images, use the finite integral Picard group
to obtain the regular comparison bimodule, and then apply Lang's
theorem to the resulting based unit-change orbit over $k$.
Throughout, each group label and the specified identity algebra
are retained.

Fix $\mathfrak R=\mathfrak R_i$, $R=R_i$, and a descent degree $c$.
First conjugate the subgroup to a pure one as above, transporting
the block and the integral corner identification.  The comparison
property applies to the conjugated data.  It passes to the basic
corners.  Explicitly, if source and target homogeneous units are
changed by $t_x$ and $t'_x$, respectively, the transport is changed
to $v\mapsto t'_xJ_x(v)t_x^{-1}$.  Covariance shows that its
multiplication rule is exactly the rule with the new crossed
factors.  The maps then restrict to the corresponding bimodule
corner.  Choose the target corner and its coordinates by applying
$\sigma^m$ to the source choices.  Thus in these fixed coordinates
the target crossed data are the actual coefficient Frobenius
images, with the same group labels.

Tensor the comparison with its successive Frobenius twists.  Both
covariance and \eqref{eq:ext:comparison-law} tensor: the two middle
factors in the product law cancel in the balanced tensor product.
After at most $c$ repetitions its exponent $n$ is divisible by
$c$.  The descent model now identifies
$\sigma^n\mathfrak R=\mathfrak R$, so its integral bimodule class
is an element $P$ of the finite group
$\mathcal P=\Pic_{\cO}(\mathfrak R)$.  Let $\theta$ be the
permutation of $\mathcal P$ induced by $\sigma^n$ and this descent
identification.  Further tensor repetitions multiply
$P,\theta(P),\theta^2(P),\ldots$.  The element $(P,\theta)$ of
the finite semidirect product
$\mathcal P\rtimes\langle\theta\rangle$ has finite order, bounded
in terms of $|\mathcal P|$.  After that many repetitions the
integral bimodule class is the identity.  The resulting total
exponent $N$ is bounded in terms of the fixed finite list and the
original comparison bound, and remains divisible by $c$.

Identify the reduced comparison bimodule with the regular
$(R,R)$-bimodule.  Put $v_x=J_x(1)$.  Covariance gives
\[
 J_x(r)=\alpha'_x(r)v_x=v_x\alpha_x(r).
\]
Since $J_x$ is invertible, multiplication by $v_x$ is invertible,
and $v_x\in R^\times$.  The product rule gives
\begin{align}
 \alpha'_x&=\operatorname{ad}(v_x)\alpha_x,
 \label{eq:ext:gauge-action}\\
 a'_{xy}&=v_x\alpha_x(v_y)a_{xy}v_{xy}^{-1}.
 \label{eq:ext:gauge-factors}
\end{align}
Thus the crossed system and its $p^N$-coefficient Frobenius image
are in the same orbit under changes of homogeneous units,
identically on $R$.

The comparison has now become a change of homogeneous units,
so it fixes the specified copy of $R$ pointwise.  To finish, we
descend this based orbit to a bounded finite field.
For fixed $R$ and $S$, all crossed systems form an affine variety
of finite type: take the automorphism matrices of $R$, the
coordinates of $a_{xy}\in R^\times$, and impose the two crossed
identities.  Invertibility can be expressed by adjoining inverse
determinants.  This variety is defined over $\F_{p^c}$.  The
unit-change group is
\[
 \mathcal U=(R^\times)^{S\setminus\{1\}},
\]
acting by \eqref{eq:ext:gauge-action}--\eqref{eq:ext:gauge-factors}.
It is connected: $R^\times$ is a nonempty open subset of the
vector space $R$.  Write $F$ for $p^N$-coefficient Frobenius.  If
$F(d)=g\cdot d$, Lang's theorem \cite{Lang1956} supplies
$h\in\mathcal U$ with $h^{-1}F(h)=g^{-1}$; hence $h\cdot d$ is
$F$-fixed.  The orbit therefore has a representative over
$\F_{p^N}$.  There are only finitely many such points for a fixed
$N$.  There are only finitely many possible bounded $N$ and
bounded-order groups $S$.  This proves based finiteness.  Transport
back through the initial conjugation gives the same conclusion
for the original restrictions.
\end{proof}

\begin{lemma}[Removing a large $p'$-kernel]\label{lem:ext:kernel-removal}
Fix an indecomposable finite-dimensional basic $k$-algebra $R$ and
a finite subgroup $\mathcal F\leq\Pic_k(R)$.  Consider crossed
algebras on $R$ by groups $T$ such that
\begin{enumerate}
\item $[T:O_{p'}(T)]$ is bounded;
\item their outer actions take values in $\mathcal F$; and
\item their restrictions to $p$-subgroups have finitely many based
isomorphism classes, for each possible abstract subgroup.
\end{enumerate}
Then their block summands have finitely many $k$-linear Morita
equivalence classes.
\end{lemma}

\begin{proof}
Let $\mathcal T$ be one of these crossed algebras and set
\[
 K=O_{p'}(T)\cap\ker(T\longrightarrow\Pic_k(R)).
\]
The index $[T:K]$ is bounded.  Change homogeneous units in the
$K$-degrees so that they centralize $R$; their factors now belong
to $Z(R)^\times$.  Since $R$ is indecomposable, its artinian center
is local, with residue field $k$, and
\begin{equation}\label{eq:ext:center-units}
 Z(R)^\times=k^\times\times U_R,
 \qquad U_R=1+\rad Z(R).
\end{equation}
This is a canonical decomposition into scalars and principal
units.  The abelian group $U_R$ has bounded $p$-power exponent:
if $(\rad Z(R))^{p^e}=0$, then $(1+z)^{p^e}=1$ for every
$z\in\rad Z(R)$.  Multiplication by $|K|$ is invertible on $U_R$,
so $H^a(K,U_R)=0$ for $a>0$.  A further central change of units
therefore makes all $K$-factors scalar.

Let $v_k$ denote these units.  Their $k$-span is a scalar twisted
group algebra $E$ of $K$, and the subalgebra on the $K$-degrees is
$R\otimes_k E$.  The algebra $E$ is semisimple by the twisted
Maschke theorem, since $p\nmid|K|$.  Every homogeneous conjugation
normalizes $E$.  Indeed it takes $v_k$ to
$z_kv_{tkt^{-1}}$, with $z_k\in Z(R)^\times$.  Comparing products
of these elements shows that the principal-unit components of
$z_k$ define a homomorphism $K\to U_R$: all original and
conjugated factors are scalar.  Such a homomorphism is trivial,
because $K$ is a $p'$-group and $U_R$ has $p$-power exponent.
Thus all $z_k$ are scalar, as required.

Write $E=\prod_j E_j$, with $E_j$ full matrix algebras over $k$,
and let $e_j$ be their identity idempotents.  Orbit sums under $T$
are central in $\mathcal T$.  In each such summand a single $e_j$
is full.  Let $T_j$ stabilize $E_j$.  Coarsening its corner grading
by $K$ gives a crossed algebra on $R\otimes E_j$ with group
\[
 Q_j=T_j/K,\qquad |Q_j|\leq[T:K].
\]
Conjugation by each homogeneous unit preserves separately $R$
and $E_j$.  By Skolem--Noether its action on $E_j$ is implemented
by a unit of $E_j$.  Dividing by that unit makes it centralize
$E_j$, without changing its action on $R$.  Its multiplication
factors now belong to the centralizer of $E_j$ in $R\otimes E_j$,
namely $R$.  Hence the corner is a matrix algebra over a crossed
algebra $\mathcal C_j$ on $R$ with grading group $Q_j$ and with
the same outer action on $R$.

The matrix factors have been removed without changing the outer
action on $R$.  This alone is insufficient for based finiteness:
we must also recover the original multiplication factors on a
Sylow subgroup.  We do so without bounding the matrix size.
Choose a Sylow $p$-subgroup $S$ of $Q_j$.  Schur--Zassenhaus in
its inverse image gives a $p$-subgroup $\widetilde S\leq T_j$
mapping isomorphically onto $S$.  Use the original homogeneous
units $u_s$ for $s\in\widetilde S$; their products have factors
$a_{s,t}\in R^\times$.  These factors centralize $E_j$, so their
conjugations on $E_j$ form an honest group action.  Choose
$c_s\in E_j^\times$ implementing it.  Then
\[
 c_sc_tc_{st}^{-1}=\gamma(s,t)\in k^\times.
\]
The corrected units $w_s=c_s^{-1}e_ju_s$ have the original actions
on $R$ and satisfy
\[
 w_sw_t=\gamma(s,t)^{-1}a_{s,t}w_{st}.
\]
Since $H^2(S,k^\times)=0$, rescaling removes $\gamma$.  Thus the
restriction of $\mathcal C_j$ to $S$ is based-isomorphic to one
of the restrictions already controlled in the hypothesis.

It remains to pass from a Sylow subgroup to a group $Q$ of bounded
order.  There are finitely many outer homomorphisms
$Q\to\mathcal F$.  Fix one and choose action representatives
$\alpha_q\in\Aut_k(R)$.  Changing homogeneous units permits
these same representatives for every crossed system with that
outer homomorphism.  If factors satisfying the crossed identities
exist, their classes under central changes of homogeneous units
form a torsor under
\[
 H^2(Q,Z(R)^\times).
\]
To see this directly, the ratio of any two factor systems with
these fixed actions is central; the associativity identity says
that this ratio is a $2$-cocycle.  A change which leaves all the
chosen actions fixed is necessarily central and gives precisely
a coboundary.

Restriction from this cohomology group to that of a Sylow
$p$-subgroup has finite kernel.  On the factor $U_R$ of
\eqref{eq:ext:center-units}, restriction followed by corestriction
is multiplication by the $p'$-index, which is invertible on this
$p$-power-exponent cohomology group.  Restriction is therefore
injective on $H^2(Q,U_R)$.  On scalars, $H^2(Q,k^\times)$ is
finite.  Indeed, for $n=|Q|$, the $n$-power map on $k^\times$ is
surjective with finite kernel $\mu_n(k)$; the long exact
cohomology sequence and annihilation by $n$ make
$H^2(Q,k^\times)$ a quotient of the finite group
$H^2(Q,\mu_n(k))$.  The decomposition
\eqref{eq:ext:center-units} is invariant under all the actions,
so these two assertions prove the finite-kernel claim.

Based Sylow finiteness gives finitely many restricted torsor
classes with the fixed actions, and each has finitely many
preimages.  There are therefore finitely many possibilities for
$\mathcal C_j$.  Each has finitely many block summands, proving
the claimed Morita finiteness.  The number or the matrix sizes
of the factors $E_j$ were never required to be bounded.
\end{proof}

\subsection{The extension criterion}

\begin{theorem}[The conditional extension criterion]
\label{thm:extension-finiteness}
Fix $p$ and $M$.  Suppose that the comparison property of
Definition~\ref{def:ext:comparison-data} holds uniformly for the
data constructed from all reduced blocks of defect order at most
$M$.  Then the blocks of all finite groups of defect order at
most $M$ have finitely many $k$-linear Morita equivalence classes.
\end{theorem}

\begin{proof}
Reduce to $B$ and construct the crossed algebra of
Lemma~\ref{lem:ext:dual-recovery}.  Its block summands include,
up to Morita equivalence, the original block.  By
Lemma~\ref{lem:ext:integral-bases} its identity-component basic
orders lie in a finite list; their reductions have a fixed finite
set of possible outer actions.  Their grading groups have a
normal $p'$-subgroup of bounded index by
Lemma~\ref{lem:ext:reduced-structure}.  Lemma~\ref{lem:ext:based-sylow}
supplies the needed based finiteness on $p$-subgroups.
Lemma~\ref{lem:ext:kernel-removal} therefore gives finitely many
Morita types for all their block summands.  Taking the finite union
over the integral base orders proves the assertion.
\end{proof}

Proposition~\ref{prop:sylow-comparison}, proved in the next section,
establishes exactly the comparison property used in
Theorem~\ref{thm:extension-finiteness}.  Combining the two results
proves the $k=\overline{\F}_p$ case of Theorem~\ref{thm:donovan};
Section~\ref{sec:coefficient-extension} transfers the resulting finite
list to the fixed algebraically closed field $K$.  The order of this exposition
introduces no extra premise: the proof of that proposition uses
the component constructions and character results, and does not
use the finite lists deduced here.  The integral finiteness theorem
has been applied only to the bounded-defect identity-component
blocks $\bar B$; the final arbitrary-group assertion is over $k$.
